\documentclass[11pt]{article}

\usepackage[margin=1in]{geometry}
\usepackage{amsmath,amssymb}
\usepackage{booktabs}
\usepackage{hyperref}

\title{Scalar-Carrier CP4 Geometry Robustness Release for HAOS-IIP}
\author{Tomislav Rupi\'c \\ Independent Researcher}
\date{April 2026}

\begin{document}
\maketitle

\begin{abstract}
This paper freezes release \texttt{51.3} of HAOS-IIP as the numbered robustness follow-up to the already-positive scalar-carrier \texttt{CP4} geometry bridge. The earlier scalar bridge had shown that one common local scalar kernel-graph family can organize Green response, heat behavior, shell-arrival structure, and first-shell low-mode organization under one common coarse Euclidean shell reconstruction. Release \texttt{51.3} asks the next honest question: does that bridge survive bounded stress on the same carrier? The answer is yes on the tested window. Across all tested mild point-set disorder cases and bounded kernel-family variants, the same scalar carrier remains compatible with the same 3D-like geometry read, while the extracted coordinate-response tensor stays close to isotropic. The result strengthens the repository's geometry authority boundary from a one-family bridge statement to a bounded robust-carrier statement, without promoting the result into broad universality, curvature, or ontology.
\end{abstract}

\tableofcontents

\section{Scope}

Release \texttt{51.3} is a narrow scalar-carrier synthesis freeze.

It does not add:
\begin{itemize}
\item a new vector validation claim beyond the already frozen \texttt{51.1} and \texttt{51.2} line;
\item a broader continuum or spacetime claim than the bounded scalar carrier actually supports;
\item broad universality across arbitrary kernels, disorder classes, or irregular substrates;
\item curvature extraction, conserved currents, or geometry-as-ontology language.
\end{itemize}

It does add:
\begin{itemize}
\item a numbered release for the first bounded robustness pass on the scalar-carrier \texttt{CP4} geometry bridge;
\item one compact authority statement covering mild point-set disorder, bounded kernel-family variation, and coordinate-response extraction on the same carrier;
\item a sharper boundary between what now counts as a robust scalar geometry read and what remains open.
\end{itemize}

\section{Where the Repository Stood Before 51.3}

Before this release, the scalar kernel-graph line had already crossed one important threshold: the same common scalar carrier had been used to organize
\begin{itemize}
\item Green response;
\item heat shell behavior;
\item shell-arrival slopes;
\item first-shell low-mode organization.
\end{itemize}

That resolved the earlier preflight mismatch in which Green response and late-phase distance-surrogate diagnostics had been living on different substrate families. But the resulting authority boundary was still narrow:

\begin{quote}
the scalar-carrier geometry bridge is positive on one tested local cubic family.
\end{quote}

Release \texttt{51.3} exists to test whether that statement survives the first bounded stress program on the same carrier rather than being left as a one-family coincidence.

\section{The 51.3 Robustness Program}

The same scalar carrier is kept fixed:
\begin{itemize}
\item scalar kernel graph on \(3D\) cubic point clouds;
\item local regime \(c_\epsilon = 0.5\);
\item same Green / heat / shell-arrival / first-shell low-mode observables;
\item same Euclidean shell reconstruction centered at the source.
\end{itemize}

The bounded robustness program then adds three controlled extensions.

\subsection{Mild point-set disorder}

The first pass perturbs the point cloud itself while preserving the same shell reconstruction logic.

\begin{itemize}
\item sizes: \(n = 11, 13\);
\item jitter fractions: \(0.00, 0.02, 0.04\);
\item kernel family fixed to \texttt{gaussian\_local}.
\end{itemize}

For jittered cases, shell bins are assigned relative to the corresponding clean-grid shell at the same \(n\), so the coarse reconstruction is not silently changed at the same time as the point set.

\subsection{Bounded kernel-family variation}

The second pass changes the local kernel law while preserving the same operator role and the same scalar carrier.

\begin{itemize}
\item sizes: \(n = 11, 13\);
\item families: \texttt{gaussian\_local}, \texttt{gaussian\_half}, \texttt{inverse\_quadratic}.
\end{itemize}

This is not a universality proof. It is only the first bounded family test.

\subsection{Coordinate-response extraction}

The third pass asks whether the same carrier still supports a metric-like response read after bounded variation. On every case the coordinate-response stiffness tensor is extracted on the orthonormalized coordinate basis:
\[
K_{\mathrm{coord}} = Q^{T} \left(-\widehat{L}_{h}\right) Q.
\]

If the same geometry read remains valid on the tested carrier, this tensor should stay close to a scalar multiple of the identity.

\section{Frozen Quantitative Summary}

Table \ref{tab:robustness_summary} records the compact bounded result frozen by \texttt{51.3}.

\begin{table}[h!]
\centering
\begin{tabular}{@{}p{3.7cm}p{2.2cm}p{8.0cm}@{}}
\toprule
Channel & Status & Frozen read \\
\midrule
Mild disorder pass & pass & All \(6/6\) tested disorder cases pass the shared geometry thresholds. \\
Kernel-family pass & pass & All \(6/6\) tested kernel-family cases pass the shared geometry thresholds. \\
Metric anisotropy & bounded & Maximum anisotropy across all tested cases is \(6.77 \times 10^{-4}\). \\
Metric off-diagonal ratio & bounded & Maximum off-diagonal ratio across all tested cases is \(1.42 \times 10^{-3}\). \\
Weakest disorder case & pass & \texttt{disorder|n=11|jitter=0.000} with shell-arrival fit \(R^2 = 0.9797\). \\
Weakest kernel case & pass & \texttt{kernel|n=13|family=gaussian\_half} with shell-arrival fit \(R^2 = 0.9732\). \\
\bottomrule
\end{tabular}
\caption{Bounded scalar-carrier robustness summary frozen by release \texttt{51.3}.}
\label{tab:robustness_summary}
\end{table}

The key point is not just that the best cases remain positive. The weakest passing cases still remain comfortably inside the bounded window.

\section{What 51.3 Now Supports}

The strongest honest scalar-geometry statement now supported by the repository is:

\begin{quote}
the scalar-carrier \texttt{CP4} geometry closure remains stable under the first bounded robustness window of mild point-set disorder, bounded kernel-family variation, and coordinate-response extraction on the tested carrier.
\end{quote}

That is stronger than the earlier bridge statement because it promotes the scalar carrier from a one-family positive example to a bounded robust carrier.

This does not replace the original bridge. It makes that bridge harder to dismiss as an artifact of one exact clean cubic setup.

\section{What 51.3 Still Does Not Say}

Release \texttt{51.3} still does not justify the following stronger claims:
\begin{itemize}
\item broad universality across large kernel families;
\item closure on strong disorder, irregular sampling, or arbitrary substrate classes;
\item curvature extraction from the same carrier;
\item conserved-current recovery;
\item a repo-wide geometry closure claim across all lines of evidence;
\item spacetime or ontology language.
\end{itemize}

Those remain later questions.

\section{Practical Consequence}

The scalar carrier is now the repository's strongest bounded geometry anchor.

It already had:
\begin{itemize}
\item the strongest bounded scalar continuum-facing control;
\item one common-family geometry bridge for Green response, heat behavior, shell-arrival structure, and low-mode organization.
\end{itemize}

Release \texttt{51.3} adds:
\begin{itemize}
\item bounded mild-disorder robustness;
\item bounded kernel-family robustness;
\item bounded coordinate-response / metric extraction robustness.
\end{itemize}

So the next honest scalar-geometry step is no longer ``does any geometry-like read survive at all?'' It is one of the following stronger programs:
\begin{itemize}
\item wider universality on the same scalar carrier;
\item stronger response / metric extraction beyond the coordinate stiffness tensor;
\item careful curvature-like or response-like closure on the same carrier.
\end{itemize}

\section{Conclusion}

Release \texttt{51.3} freezes the bounded robustness follow-up to the scalar-carrier \texttt{CP4} geometry bridge.

The result is positive on the tested window. Mild point-set disorder, bounded kernel-family variation, and coordinate-response extraction all remain compatible with the same 3D-like geometry read on the same scalar carrier. That upgrades the scalar geometry line from a one-family bridge result to a bounded robust-carrier statement.

That is the right stopping point for this release: stronger than the initial bridge, but still disciplined about what the repository has not yet earned.

\end{document}
